\section{Limitations}
\label{sec:limits}

The value of an inventory of this kind depends on being candid about what it
is not, so the limitations are set out at length rather than compressed.

\textbf{The search was targeted, not systematic.} No registered search
strategy was followed, no databases were searched exhaustively, and no
screening log was kept. Relevant measurement studies may exist that were not
located, and a step graded E5 here means no evidence was found rather than
that none exists. A systematic search is the single largest available upgrade
to this work.

\textbf{The scores are structured expert judgement, not measurement.} The
instrument makes judgement explicit, anchored and arguable, which is a real
gain over an unstated judgement. It does not make it empirical. A reader who
disagrees with a score can see which anchor was applied and argue with that
specific item, and that is the intended use.

\textbf{Inter-rater agreement has not been measured.} Every score in this
brief reflects one rater's reading of the anchors. A protocol is registered in
advance and the analysis code is written and self-tested, but no second rater
has yet scored anything. Until that study runs, the reproducibility of these
scores is unknown, and they should be read as a documented position rather
than as a measurement.

\textbf{Negative measurement results are weaker than they appear.} The studies
that recovered nothing used model organisms at chosen concentrations, in single
laboratories, under correct technique, with containment in place, and with
finite detection limits. ``Nothing was recovered above 4~CFU\,m$^{-3}$'' is not
``nothing is released'', and for a low-dose inhalational organism the residual
is not necessarily negligible.

\textbf{The cross-study comparison in Figure~\ref{fig:measured} is
illustrative.} The studies differ in organism, sampler, flow rate, room,
ventilation and specimen matrix, and were not designed for comparison. The
figure is offered as a map of coverage, not as a calibrated ranking.

\textbf{Baseline control scores characterise reported practice, not surveyed
practice.} No survey of what laboratories actually do was conducted. The
baseline column describes common practice as described in the literature and
guidance, which is a characterisation rather than an estimate.

\textbf{Agent profiles are stylised.} Three profiles stand in for a wide
diversity of organisms. A laboratory with a specific organism in view should
score the consequence dimension for that organism rather than take a profile
off the shelf.

\textbf{The inventory describes US non-waived practice.} Equipment, staffing
patterns and regulation differ elsewhere. The step definitions should transfer;
the baseline control scores very likely do not.

\textbf{One control recommendation is more cautious than the evidence may
warrant.} The claim about matrix-only inactivation on a mass spectrometry
target was located only at second hand and is excluded, so the MALDI-TOF
control set is drawn tighter than a full reading might justify. For the same
reason, no standard designation or certification interval is named for
biological safety cabinets; the inventory instead requires each laboratory to
record the standard and interval it applies.

\textbf{No numeric wait interval is given for opening a centrifuge.} No
interval in the located evidence derives from measurement of a clinical
centrifuge, and a number stated here would acquire unwarranted authority.

\textbf{The survey items are unvalidated.} They are candidates derived from
evidence gaps, not a tested instrument.

\textbf{ISO 35001:2019 is used structurally.} Clause numbers and titles index
the elements of the inventory. No text of the standard is reproduced, and
nothing here asserts or confers conformity.

\section{Conclusion}

Seven bench steps make up most of the physical handling in routine non-waived
diagnostic testing. This brief defines them, states what is known about aerosol
release from each, grades the strength of that knowledge, and records the
controls the evidence supports.

The evidence is thinner than the volume of guidance implies, and it is
distributed opposite to the hazard. Three of seven steps carry direct
measurement. The two steps with the highest release scores and the only
documented human infections carry none. The highest-throughput step in the
laboratory carries no evidence of any kind.

For tuberculosis work the measured record is both unusually good and unusually
narrow. Cartridge-based molecular testing produced no culturable bioaerosol at
any stage, against 16 and 325~CFU\,m$^{-3}$ for the smear microscopy it
displaces, which is as clean a result as biosafety evidence offers. It covers
one step of six on the pathway. The other five are governed by incident
evidence or by nothing.

Four measurement studies would resolve most of the uncertainty documented here:
one high-volume automated analyser covering normal running and cover-open
maintenance separately; open-bench subculture and plate manipulation; MALDI-TOF
target preparation; and the pre-analytical handling of respiratory specimens
before rule-out. None is technically difficult. Each would convert a row of
this inventory from precaution to evidence.

Until then, a laboratory assessing its own bench should know which of its
controls rest on measurement and which rest on custom. Making that distinction
visible is what this brief is for.

\section*{Data and code availability}

The inventory, the extracted evidence with sources and access dates, the
grading, the scoring instrument, the figure and analysis code, and the
provenance log are openly available at
\url{https://github.com/DrAdedamola/clinical-lab-aerosol-biorisk-reference}.
Instrument, data and documents are released under CC~BY~4.0; code under the
MIT licence. Every figure in this brief is generated from the project's own
data files by the included scripts.

\section*{Competing interests}

None declared.

\section*{Funding}

None.

\section*{Use of artificial intelligence}

AI assistance was used for code, data wrangling, literature triage and drafting
mechanics. Every analytic decision, every score and every citation is the
author's and was checked against its source. Errors identified during the
author's verification pass, including a scoring defect in the aggregation of a
reverse-anchored item and four citation corrections, were corrected before
release and are recorded in the project release notes.
